\chapter{Rappresentazione delle informazioni}
\label{chap:rappresentazione-informazioni}

Come anticipato nella \cref{sec:fi-informatica-calcolo}, un elaboratore
memorizza l'informazione mediante sequenze di bit. L'interpretazione
di una sequenza dipende dalla codifica adottata: gli stessi bit possono
rappresentare numeri, caratteri oppure altri dati.\footnote{Riferimento didattico:
S. Canazza, \emph{Rappresentazione delle informazioni},
\nolinkurl{2026_FI_01_canazza_RappresentazioneInformazioni.pdf}.}

\section{Sistemi di numerazione posizionali}
\label{sec:fi-notazione-posizionale}

Nei sistemi di numerazione \emph{posizionali} il valore di una cifra dipende
dalla sua posizione e dalla base. Il sistema decimale impiega le cifre da
\(0\) a \(9\); quello binario le cifre \(0\) e \(1\).

\begin{definition}[Valore posizionale]
\label{def:fi-valore-posizionale}
Siano \(b\ge2\) una base intera e \(a_k\in\{0,\ldots,b-1\}\) le cifre di
una rappresentazione con \(n\ge1\) cifre intere e \(m\ge0\) cifre
frazionarie. Il valore della scrittura
\[
  (a_{n-1}\cdots a_0.a_{-1}\cdots a_{-m})_b
\]
è
\[
  x=\sum_{k=0}^{n-1}a_kb^k+\sum_{j=1}^{m}a_{-j}b^{-j}.
\]
Per \(m=0\) la seconda somma è nulla e il separatore si omette.
\end{definition}

Si utilizza il punto come separatore della parte frazionaria e il pedice
per indicare la base.

\begin{example}[Pesi delle cifre decimali]
\begin{align*}
  (753)_{10}&=7\cdot10^2+5\cdot10+3=753,\\
  (2048)_{10}&=2\cdot10^3+4\cdot10+8=2048,\\
  (4.34)_{10}&=4+3\cdot10^{-1}+4\cdot10^{-2}.
\end{align*}
\end{example}

\section{Notazione binaria e conversioni}
\label{sec:fi-conversioni-binarie}

Il \emph{bit} (\emph{binary digit}) è una cifra binaria. La base due è
adatta ai circuiti digitali, nei quali è possibile distinguere due stati
logici. La conversione da binario a decimale si ottiene applicando la
\cref{def:fi-valore-posizionale}.

\begin{example}[Da binario a decimale]
\begin{align*}
  (1101)_2&=2^3+2^2+1=13,\\
  (1.101)_2&=1+2^{-1}+2^{-3}=1.625.
\end{align*}
\end{example}

Viceversa, per convertire un numero decimale non negativo in binario
si trattano separatamente la parte intera e quella frazionaria.

\subsection{Divisioni successive per la parte intera}

Sia \(N\in\N\), con la convenzione \(0\in\N\). Per \(N>0\), si divide
ripetutamente per due, ottenendo quozienti \(q_i\) e resti \(r_i\):
\[
  q_0=N,\qquad q_i=2q_{i+1}+r_i,
  \qquad r_i\in\{0,1\}.
\]
Quando il quoziente diventa zero, la sequenza binaria è formata dai resti
letti \emph{in ordine inverso}. Infatti, sostituendo successivamente i
quozienti si ricava
\[
  N=r_0+2r_1+\cdots+2^{t-1}r_{t-1}.
\]
Per \(N=0\) la rappresentazione è \(0\).

\begin{example}[Conversione di cento]
\begin{alignat*}{3}
  100&=2\cdot50+0, &\quad 50&=2\cdot25+0, &\quad25&=2\cdot12+1,\\
  12&=2\cdot6+0, & 6&=2\cdot3+0, & 3&=2\cdot1+1,\\
  1&=2\cdot0+1.
\end{alignat*}
Leggendo i resti dal basso verso l'alto,
\[
  (100)_{10}=(1100100)_2=2^6+2^5+2^2.
\]
\end{example}

\subsection{Moltiplicazioni successive per la parte frazionaria}

Sia \(f_0\in[0,1)\). Si moltiplica ripetutamente la frazione per due:
\[
  2f_{j-1}=d_j+f_j,
  \qquad d_j\in\{0,1\},\quad f_j\in[0,1).
\]
Le parti intere \(d_j\) forniscono i bit frazionari, letti
\emph{nell'ordine di produzione}. Dopo \(t\) passi,
\[
  f_0=\sum_{j=1}^{t}d_j2^{-j}+f_t2^{-t}.
\]
La conversione termina se \(f_t=0\); se una frazione residua si ripete,
la sequenza di bit è periodica da quel punto.

\begin{example}[Frazione binaria finita]
\[
  0.625\cdot2=1+0.25,\qquad
  0.25\cdot2=0+0.5,\qquad
  0.5\cdot2=1.
\]
Pertanto \((0.625)_{10}=(0.101)_2\).
\end{example}

\begin{example}[Frazione binaria periodica]
\label{ex:fi-frazione-periodica}
Per \(0.35=7/20\) si ottiene
\begin{align*}
  2\cdot0.35&=0+0.70, & 2\cdot0.70&=1+0.40,\\
  2\cdot0.40&=0+0.80, & 2\cdot0.80&=1+0.60,\\
  2\cdot0.60&=1+0.20, & 2\cdot0.20&=0+0.40.
\end{align*}
Poiché \(0.40\) ricompare, il periodo è \(0110\). Quindi
\[
  (0.35)_{10}=(0.01\overline{0110})_2,
  \qquad
  (100.35)_{10}=(1100100.01\overline{0110})_2.
\]
\end{example}

\begin{proposition}[Frazioni con rappresentazione binaria finita]
\label{prop:fi-frazione-binaria-finita}
Una frazione \(p/q\in[0,1)\), ridotta ai minimi termini e con \(q>0\),
ammette una rappresentazione binaria finita se e solo se \(q\) è una
potenza di due.
\end{proposition}

\begin{proof}
Una scrittura con \(m\) cifre frazionarie ha valore \(K/2^m\), quindi
il denominatore ridotto divide \(2^m\). Viceversa, se \(q=2^h\), bastano
al più \(h\) cifre. Negli altri casi i resti della conversione assumono
soltanto un numero finito di valori non nulli: uno si ripete e la
rappresentazione è periodica da una certa posizione.
\end{proof}

\section{Numeri senza segno e modulo e segno}
\label{sec:fi-interi-rappresentazione}

Fissata una lunghezza di \(n\ge1\) bit, si hanno \(2^n\) configurazioni.
Nell'interpretazione \emph{senza segno}, la sequenza \(a_{n-1}\cdots a_0\) vale
\[
  U=\sum_{k=0}^{n-1}a_k2^k,
  \qquad U\in[0,2^n-1]\cap\Z.
\]
Con otto bit si rappresentano quindi gli interi da \(0\) a \(255\).
Il bit più a sinistra è detto \emph{più significativo} (MSB); quello
più a destra \emph{meno significativo} (LSB).

\begin{definition}[Modulo e segno]
Per \(n\ge2\), nella rappresentazione con \emph{modulo e segno} il bit
più significativo \(s\) indica il segno, mentre gli altri \(n-1\) bit
rappresentano il modulo \(M\):
\[
  a=(-1)^sM,\qquad
  M=\sum_{k=0}^{n-2}a_k2^k.
\]
\end{definition}

Le sequenze \(001100\) e \(101100\) su sei bit rappresentano
rispettivamente \(+12\) e \(-12\). L'intervallo dei valori è
\[
  [-(2^{n-1}-1),\,2^{n-1}-1]\cap\Z.
\]
Le configurazioni sono \(2^n\), ma i valori distinti sono \(2^n-1\):
\(00\cdots0\) e \(10\cdots0\) rappresentano entrambe zero.
Inoltre, la somma richiede di distinguere operandi con segno uguale,
per i quali si sommano i moduli, e con segno opposto, per i quali
si confrontano e sottraggono i moduli.

\section{Complemento a due}
\label{sec:fi-complemento-due}

Il \emph{complemento a due} evita la doppia rappresentazione dello zero
e permette di eseguire le somme mediante l'addizione binaria.

\begin{definition}[Codifica in complemento a due]
\label{def:fi-complemento-due}
Sia \(n\ge2\). Per ogni intero
\[
  a\in[-2^{n-1},2^{n-1}-1]\cap\Z,
\]
la codifica in complemento a due è la rappresentazione senza segno
su \(n\) bit di
\[
  C_n(a)=
  \begin{cases}
    a, & a\ge0,\\
    a+2^n, & a<0.
  \end{cases}
\]
\end{definition}

Per decodificare una sequenza di valore senza segno \(U\), si pone
\[
  a=\begin{cases}
    U, & U<2^{n-1},\\
    U-2^n, & U\ge2^{n-1}.
  \end{cases}
\]
Equivalentemente, il bit più significativo ha peso negativo:
\begin{equation}
  a=-a_{n-1}2^{n-1}+\sum_{k=0}^{n-2}a_k2^k.
  \label{eq:fi-pesi-complemento-due}
\end{equation}
Il bit più significativo indica quindi il segno, ma gli altri bit di
un numero negativo \emph{non} ne rappresentano il modulo. Gli estremi
sono \(-2^{n-1}\) e \(2^{n-1}-1\); su otto bit si ha
\([-128,127]\cap\Z\) e su trentadue bit
\([-2147483648,2147483647]\cap\Z\), intervallo del tipo
\texttt{int} di Java.

\subsection{Codifica, inversione e decodifica}

Per ottenere la codifica di \(-m\) su \(n\) bit, con
\(1\le m\le2^{n-1}\), si calcola \(2^n-m\). Equivalentemente, si
scrive \(m\) su \(n\) bit, si invertono tutti i bit e si aggiunge uno.
Infatti, l'inversione produce \(2^n-1-m\).

\begin{example}[Un intero negativo su otto bit]
Per \(-13\) si ha
\[
  2^8-13=243=(1111\,0011)_2.
\]
Anche l'inversione con incremento fornisce
\[
  0000\,1101\longrightarrow1111\,0010
  \xrightarrow{+1}1111\,0011.
\]
Dall'\cref{eq:fi-pesi-complemento-due} segue
\[
  -128+64+32+16+2+1=-13.
\]
Analogamente, \(-14\) corrisponde a \(1111\,0010\).
\end{example}

\begin{warning}[Il formato determina il valore]
\label{warn:fi-larghezza-interi}
La sequenza \(1111\,1111\) rappresenta \(255\) senza segno, \(-127\)
con modulo e segno e \(-1\) in complemento a due su otto bit. In
complemento a due su sedici bit, \(0000\,0000\,1111\,1111\) vale
invece \(+255\). È necessario specificare formato e larghezza.
\end{warning}

Nell'\emph{estensione del segno} si replica il bit più significativo
per aumentare la larghezza senza cambiare il valore. Per esempio,
\(1111\,0011\) su otto bit rappresenta \(-13\); la sua estensione
a sedici bit è
\[
  (1111\,1111\,1111\,0011)_2.
\]
La negazione di \(-2^{n-1}\) non è rappresentabile con \(n\)
bit in complemento a due.

\section{Aritmetica intera e overflow}
\label{sec:fi-aritmetica-overflow}

L'addizione binaria conserva \(n\) bit e scarta il riporto finale:
il risultato è la somma delle codifiche \emph{modulo} \(2^n\).
Nella colonna \(k\), con riporto entrante \(c_k\), vale
\[
  a_k+b_k+c_k=r_k+2c_{k+1},\qquad c_0=0.
\]
Se la somma matematica non appartiene all'intervallo rappresentabile,
si verifica \emph{overflow} (trabocco).

\begin{example}[Riporto scartato senza overflow]
Su otto bit,
\[
  \begin{array}{r@{\;}r}
      &0000\,0101\\
    + &1111\,1110\\ \hline
    1 &0000\,0011
  \end{array}
\]
rappresenta \(5+(-2)=3\). Il riporto finale viene scartato, ma
non si verifica overflow con segno.
\end{example}

\begin{proposition}[Criterio di overflow nell'addizione]
\label{prop:fi-overflow-somma}
Nella somma di due interi su \(n\) bit in complemento a due,
si verifica overflow se e solo se gli operandi hanno lo stesso segno
e il risultato ha segno opposto. Equivalentemente, il riporto
entrante nel bit di segno \(c_{n-1}\) è diverso dal riporto uscente
\(c_n\).
\end{proposition}

\begin{proof}
La somma di interi con segno opposto è compresa fra gli operandi e non
può eccedere l'intervallo disponibile. La somma di due non negativi
può superare solo il massimo e produrre un primo bit uguale a uno;
quella di due negativi può scendere sotto il minimo e produrre un
primo bit nullo. Applicando la regola dell'addizione alla colonna di
segno, i due casi corrispondono rispettivamente a
\((c_{n-1},c_n)=(1,0)\) e \((0,1)\).
\end{proof}

\begin{example}[Overflow positivo e negativo su quattro bit]
L'intervallo è \([-8,7]\cap\Z\). Si ha
\[
  0110+0010=1000,\qquad 1101+1001=1\,0110.
\]
Nel primo caso \(6+2=8\) produce \(-8\); nel secondo
\((-3)+(-7)=-10\) produce \(+6\) dopo lo scarto del riporto.
Entrambi i risultati sono corretti modulo \(16\), ma non come
somme di interi con segno.
\end{example}

La sottrazione \(a-b\) si esegue sommando ad \(a\) il complemento a
due di \(b\). Si verifica overflow quando gli operandi della
sottrazione hanno segni opposti e il risultato ha segno diverso da
quello di \(a\). Il riporto finale, da solo, non determina l'overflow
con segno (\cref{prop:fi-overflow-somma}).

\section{Virgola fissa e virgola mobile}
\label{sec:fi-virgola-mobile}

Con \(n\) bit si distinguono soltanto \(2^n\) configurazioni, quindi
non tutti i numeri reali sono rappresentabili esattamente. Alcune
frazioni, tuttavia, hanno una scrittura binaria finita, come mostrato
nella \cref{sec:fi-conversioni-binarie}.

Nella rappresentazione in \emph{virgola fissa} la posizione del
separatore è prestabilita. Usando \(n\) bit senza segno, di cui
\(m\) frazionari, si rappresentano valori \(U2^{-m}\), con
\(0\le U\le2^n-1\), distanziati di \(2^{-m}\).

Nella rappresentazione in \emph{virgola mobile} un numero finito non
nullo è espresso nella forma
\[
  x=(-1)^sM\,b^E,
\]
dove \(s\) indica il segno, \(M\) è il \emph{significando} (o
\emph{mantissa}) ed \(E\) l'esponente. La \emph{normalizzazione}
richiede \(1\le M<b\). In base due, per esempio,
\[
  (10.5)_{10}=(1010.1)_2=(1.0101)_2\cdot2^3.
\]
Il numero di cifre significative determina la precisione, mentre
l'esponente determina l'ordine di grandezza.

\subsection{I formati IEEE 754 binary32 e binary64}
\label{sec:fi-ieee754}

Lo standard IEEE 754 definisce, fra gli altri, i formati \emph{binary32}
e \emph{binary64}, rispettivamente a singola e doppia precisione.\footnote{
Oracle, \emph{Numerical Computation Guide},
\href{https://docs.oracle.com/cd/E19957-01/806-3568/ncg_math.html}{capitolo 2, IEEE Arithmetic}.}
Ciascuno comprende un bit di segno \(s\), un campo per l'esponente e
uno per la frazione. I parametri sono riportati nella
\cref{tab:fi-formati-ieee}.

\begin{table}[htbp]
  \centering
  \caption{Parametri dei formati IEEE 754 binari.}
  \label{tab:fi-formati-ieee}
  \begin{tabular}{lcc}
    \toprule
    Parametro & binary32 & binary64 \\
    \midrule
    Bit complessivi & 32 & 64 \\
    Segno & 1 & 1 \\
    Esponente \(w\) & 8 & 11 \\
    Frazione \(t\) & 23 & 52 \\
    Bias \(B\) & 127 & 1023 \\
    Bit significativi normali & 24 & 53 \\
    Esponenti normali & \(-126\) a \(127\) & \(-1022\) a \(1023\) \\
    \bottomrule
  \end{tabular}
\end{table}

Sia \(e\) il valore intero senza segno del campo esponente e \(F\)
quello del campo frazione. Il \emph{bias} \(B=2^{w-1}-1\) permette di
codificare gli esponenti senza usare il complemento a due: per i
numeri normali l'esponente effettivo è \(E=e-B\). I campi
\(e=0\) ed \(e=2^w-1\) sono riservati.

\begin{definition}[Valori normali e subnormali]
\label{def:fi-valori-ieee}
Nei formati della \cref{tab:fi-formati-ieee}, per
\(1\le e\le2^w-2\) il valore \emph{normale} è
\[
  x=(-1)^s\left(1+\frac{F}{2^t}\right)2^{e-B}.
\]
Per \(e=0\) e \(F\ne0\), il valore \emph{subnormale} è
\[
  x=(-1)^s\frac{F}{2^t}\,2^{1-B}.
\]
\end{definition}

Il bit iniziale \(1\) del significando normale è implicito;
nei subnormali il bit iniziale è invece \(0\).

\begin{example}[Codifica di un valore normale]
Il numero \(10.5=(1.0101)_2\cdot2^3\), codificato in binary32,
ha \(s=0\), \(e=3+127=130=(10000010)_2\) e frazione
\(0101\) seguita da diciannove zeri:
\[
  \underbrace{0}_{s}\quad
  \underbrace{10000010}_{e}\quad
  \underbrace{01010000000000000000000}_{\text{frazione}}.
\]
La decodifica restituisce \((1+2^{-2}+2^{-4})2^3=10.5\).
\end{example}

\subsection{Valori limite, zeri e configurazioni speciali}

Per i formati della \cref{tab:fi-formati-ieee}, il minimo normale
positivo, il minimo subnormale positivo e il massimo finito positivo
sono rispettivamente
\[
  2^{1-B},\qquad
  2^{1-B-t},\qquad
  (2-2^{-t})2^{2^w-2-B}.
\]
In binary32 valgono
\begin{align*}
  x_{\min,\mathrm{norm}}&=2^{-126}\approx1.1755\cdot10^{-38},\\
  x_{\min,\mathrm{sub}}&=2^{-149}\approx1.4013\cdot10^{-45},\\
  x_{\max}&=(2-2^{-23})2^{127}\approx3.4028\cdot10^{38}.
\end{align*}
In binary64 i corrispondenti estremi sono \(2^{-1022}\),
\(2^{-1074}\) e \((2-2^{-52})2^{1023}\).
I subnormali consentono un \emph{underflow graduale} verso zero,
con una precisione relativa decrescente.

Le configurazioni riservate sono riepilogate nella
\cref{tab:fi-valori-speciali}.

\begin{table}[htbp]
  \centering
  \caption{Interpretazione dei campi IEEE 754.}
  \label{tab:fi-valori-speciali}
  \begin{tabular}{lll}
    \toprule
    Esponente & Frazione & Valore \\
    \midrule
    \(e=0\) & \(F=0\) & \(+0\) oppure \(-0\) \\
    \(e=0\) & \(F\ne0\) & Subnormale \\
    \(1\le e\le2^w-2\) & Qualsiasi \(F\) & Normale \\
    \(e=2^w-1\) & \(F=0\) & \(+\infty\) oppure \(-\infty\) \\
    \(e=2^w-1\) & \(F\ne0\) & NaN \\
    \bottomrule
  \end{tabular}
\end{table}

\emph{NaN} (\emph{Not a Number}) rappresenta risultati non numerici,
come \(0/0\) e \(+\infty-(+\infty)\) nell'aritmetica IEEE 754.
Un confronto di uguaglianza con NaN è falso anche quando si confronta
NaN con sé stesso. Il comportamento delle operazioni e la segnalazione
degli errori dipendono anche dal linguaggio e dall'ambiente: in Python,
per esempio, \texttt{1.0/0.0} solleva \texttt{ZeroDivisionError}.

\section{Precisione e arrotondamento}
\label{sec:fi-precisione-arrotondamento}

Per i numeri normali positivi nell'intervallo
\([2^E,2^{E+1})\), la distanza fra valori consecutivi è
\[
  \delta_E=2^{E-t},
\]
dove \(t\) è il numero di bit della frazione (\cref{tab:fi-formati-ieee}).
In binary32 il passo è \(1\) in corrispondenza di \(2^{23}\) e
\(16\) in corrispondenza di \(2^{27}\). Nei subnormali il passo
resta \(2^{1-B-t}\).

Se un risultato non è rappresentabile esattamente, viene
\emph{arrotondato}. La modalità predefinita IEEE 754 è
\emph{al più vicino, con parità in caso di equidistanza}
(\emph{ties to even}); sono previste anche le modalità verso zero,
verso \(+\infty\) e verso \(-\infty\). L'aritmetica a precisione
finita distingue l'errore di rappresentazione degli operandi da quello
di arrotondamento del risultato.

\begin{example}[Il valore decimale \(4.35\)]
\label{ex:fi-arrotondamento-prodotto}
Poiché \(4.35=87/20\), il valore non ha una rappresentazione binaria
finita, per la \cref{prop:fi-frazione-binaria-finita}.
In binary64, arrotondato al più vicino, vale
\[
  \widehat{x}\approx4.3499999999999996.
\]
La moltiplicazione, arrotondata nello stesso formato, produce
\[
  \operatorname{fl}(\widehat{x}\cdot100)
  =435-2^{-44}\approx434.99999999999994,
\]
anziché \(435\). In binary32, invece, il prodotto arrotondato
è \(435\): l'esito dipende dal formato.
\end{example}

\begin{example}[Allineamento degli esponenti e perdita di cifre]
\begin{align*}
  10.5+0.125
    &=(1.0101)_2\cdot2^3+(1.0)_2\cdot2^{-3}\\
    &=\bigl((1.0101)_2+(0.000001)_2\bigr)2^3\\
    &=(1.010101)_2\cdot2^3=10.625.
\end{align*}
Il risultato richiede sette bit significativi; in binary32 e binary64
è esatto. In un formato didattico con cinque bit significativi e
arrotondamento al più vicino, il passo presso \(10.5\) è \(0.5\):
la somma viene arrotondata a \(10.5\).
\end{example}

Un addendo piccolo può dunque non modificare il valore memorizzato.
Inoltre, l'addizione in virgola mobile non è, in generale, associativa:
le approssimazioni possono variare con l'ordine delle operazioni.

\section{Notazione esadecimale e ottale}
\label{sec:fi-esadecimale-ottale}

Per scrivere più concisamente le sequenze binarie si usa anche la
base sedici, che impiega le cifre
\(0,\ldots,9,\mathrm{A},\ldots,\mathrm{F}\),
con \(\mathrm{A}=10,\ldots,\mathrm{F}=15\). Poiché \(16=2^4\),
una cifra esadecimale corrisponde a quattro bit. I gruppi si formano
partendo dal separatore: verso sinistra per la parte intera, verso
destra per quella frazionaria.

\begin{example}[Raggruppamento in base sedici]
\begin{align*}
  (0111\,1111)_2&=(\mathrm{7F})_{16}=127,\\
  (1011\,0110)_2&=(\mathrm{B6})_{16},\\
  (1101.101)_2&=(\mathrm{D.A})_{16}=13.625.
\end{align*}
Nell'ultimo esempio si aggiunge uno zero \emph{a destra} della
parte frazionaria. La notazione \texttt{0x7F} indica comunemente
\((\mathrm{7F})_{16}\).
\end{example}

La base otto impiega le cifre da \(0\) a \(7\). Poiché \(8=2^3\),
si raggruppano i bit tre a tre:
\begin{align*}
  (100010)_2&=(100\,010)_2=(42)_8=34,\\
  (11100010)_2&=(011\,100\,010)_2=(342)_8=226,\\
  (31)_8&=3\cdot8+1=25.
\end{align*}
Nelle conversioni inverse si sostituisce ogni cifra esadecimale con
quattro bit e ogni cifra ottale con tre bit. La base non determina
il segno:
per interpretare una sequenza come intero con segno occorre conoscere
il formato (\cref{warn:fi-larghezza-interi}).

\section{Rappresentazione dei caratteri}
\label{sec:fi-caratteri}

Per rappresentare un testo si associano codici numerici ai suoi
caratteri. Il \emph{punto di codice} identifica un carattere nel
repertorio; il \emph{formato di codifica} stabilisce come quel valore
viene memorizzato in byte. Uno stesso punto di codice può quindi avere
rappresentazioni in byte diverse.

\subsection{ASCII e codifiche a otto bit}

Il codice \emph{ASCII} (\emph{American Standard Code for Information
Interchange}) usa sette bit e comprende \(128\) codici, da \(0\) a
\(127\). I codici \(0\)--\(31\) e \(127\) sono di controllo;
i codici \(32\)--\(126\) sono stampabili.

\begin{table}[htbp]
  \centering
  \caption{Codici e intervalli ASCII principali.}
  \label{tab:fi-ascii}
  \begin{tabular}{lll}
    \toprule
    Codice & Simbolo & Significato \\
    \midrule
    9 & \texttt{\textbackslash t} & Tabulazione orizzontale (HT) \\
    10 & \texttt{\textbackslash n} & Nuova riga (LF) \\
    13 & \texttt{\textbackslash r} & Ritorno a inizio riga (CR) \\
    32 & Spazio & Spazio \\
    48--57 & \texttt{0}--\texttt{9} & Cifre decimali \\
    65--90 & \texttt{A}--\texttt{Z} & Lettere maiuscole \\
    97--122 & \texttt{a}--\texttt{z} & Lettere minuscole \\
    127 & \texttt{DEL} & Carattere di controllo \\
    \bottomrule
  \end{tabular}
\end{table}

I codici di \texttt{0} e \texttt{1} sono rispettivamente \(48\) e
\(49\), distinti dai valori numerici zero e uno. Le sequenze
\texttt{\textbackslash n}, \texttt{\textbackslash r} e
\texttt{\textbackslash t} sono notazioni impiegate nei linguaggi di
programmazione per indicare caratteri di controllo.

ASCII non comprende le lettere accentate. La denominazione
\emph{ASCII esteso} indica diverse codifiche a otto bit, non uno
standard unico: ISO-8859-1 e Windows-1252, per esempio, condividono
i primi \(128\) codici, ma assegnano significati differenti ad alcuni
byte successivi.

\subsection{Unicode e formati UTF}

\emph{Unicode} assegna ai caratteri punti di codice indicati nella
forma \texttt{U+XXXX}. Lo spazio dei punti di codice arriva a
\texttt{U+10FFFF}; i primi \(128\) corrispondono ad ASCII.
Unicode non prescrive due byte per carattere e non tutti i punti
di codice sono assegnati.\footnote{Unicode Consortium,
\emph{The Unicode Standard},
\href{https://www.unicode.org/versions/Unicode17.0.0/core-spec/chapter-2/}{capitolo 2}.}
L'intervallo \texttt{U+D800}--\texttt{U+DFFF} è riservato ai surrogati.

I formati principali sono \emph{UTF-8}, che impiega da uno a quattro
byte, \emph{UTF-16}, che impiega una o due unità di sedici bit, e
\emph{UTF-32}, che impiega un'unità di trentadue bit. In UTF-8
le sequenze ASCII restano invariate. UTF-16 utilizza coppie di
surrogati per i punti di codice superiori a \texttt{U+FFFF};
per la memorizzazione delle unità su più byte occorre specificare
l'ordine dei byte.

\begin{example}[Punto di codice e byte memorizzati]
La lettera \texttt{A}, punto di codice \texttt{U+0041}, occupa
il byte \texttt{41} in UTF-8. La lettera \texttt{è}, punto di
codice \texttt{U+00E8}, occupa \texttt{C3 A8} in UTF-8,
\texttt{00 E8} in UTF-16BE e \texttt{E8 00} in UTF-16LE.
In ISO-8859-1 occupa invece il byte \texttt{E8}.
\end{example}

Un carattere visualizzato può comprendere più punti di codice:
numero di caratteri, punti di codice e byte non coincidono
necessariamente.
